% Options for packages loaded elsewhere
% Options for packages loaded elsewhere
\PassOptionsToPackage{unicode}{hyperref}
\PassOptionsToPackage{hyphens}{url}
\PassOptionsToPackage{dvipsnames,svgnames,x11names}{xcolor}
%
\documentclass[
  english,
  a4paper,
]{article}
\usepackage{xcolor}
\usepackage{amsmath,amssymb}
\setcounter{secnumdepth}{5}
\usepackage{iftex}
\ifPDFTeX
  \usepackage[T1]{fontenc}
  \usepackage[utf8]{inputenc}
  \usepackage{textcomp} % provide euro and other symbols
\else % if luatex or xetex
  \usepackage{unicode-math} % this also loads fontspec
  \defaultfontfeatures{Scale=MatchLowercase}
  \defaultfontfeatures[\rmfamily]{Ligatures=TeX,Scale=1}
\fi
\usepackage{lmodern}
\ifPDFTeX\else
  % xetex/luatex font selection
\fi
% Use upquote if available, for straight quotes in verbatim environments
\IfFileExists{upquote.sty}{\usepackage{upquote}}{}
\IfFileExists{microtype.sty}{% use microtype if available
  \usepackage[]{microtype}
  \UseMicrotypeSet[protrusion]{basicmath} % disable protrusion for tt fonts
}{}
\makeatletter
\@ifundefined{KOMAClassName}{% if non-KOMA class
  \IfFileExists{parskip.sty}{%
    \usepackage{parskip}
  }{% else
    \setlength{\parindent}{0pt}
    \setlength{\parskip}{6pt plus 2pt minus 1pt}}
}{% if KOMA class
  \KOMAoptions{parskip=half}}
\makeatother
% Make \paragraph and \subparagraph free-standing
\makeatletter
\ifx\paragraph\undefined\else
  \let\oldparagraph\paragraph
  \renewcommand{\paragraph}{
    \@ifstar
      \xxxParagraphStar
      \xxxParagraphNoStar
  }
  \newcommand{\xxxParagraphStar}[1]{\oldparagraph*{#1}\mbox{}}
  \newcommand{\xxxParagraphNoStar}[1]{\oldparagraph{#1}\mbox{}}
\fi
\ifx\subparagraph\undefined\else
  \let\oldsubparagraph\subparagraph
  \renewcommand{\subparagraph}{
    \@ifstar
      \xxxSubParagraphStar
      \xxxSubParagraphNoStar
  }
  \newcommand{\xxxSubParagraphStar}[1]{\oldsubparagraph*{#1}\mbox{}}
  \newcommand{\xxxSubParagraphNoStar}[1]{\oldsubparagraph{#1}\mbox{}}
\fi
\makeatother


\usepackage{longtable,booktabs,array}
\newcounter{none} % for unnumbered tables
\usepackage{calc} % for calculating minipage widths
% Correct order of tables after \paragraph or \subparagraph
\usepackage{etoolbox}
\makeatletter
\patchcmd\longtable{\par}{\if@noskipsec\mbox{}\fi\par}{}{}
\makeatother
% Allow footnotes in longtable head/foot
\IfFileExists{footnotehyper.sty}{\usepackage{footnotehyper}}{\usepackage{footnote}}
\makesavenoteenv{longtable}
\usepackage{graphicx}
\makeatletter
\newsavebox\pandoc@box
\newcommand*\pandocbounded[1]{% scales image to fit in text height/width
  \sbox\pandoc@box{#1}%
  \Gscale@div\@tempa{\textheight}{\dimexpr\ht\pandoc@box+\dp\pandoc@box\relax}%
  \Gscale@div\@tempb{\linewidth}{\wd\pandoc@box}%
  \ifdim\@tempb\p@<\@tempa\p@\let\@tempa\@tempb\fi% select the smaller of both
  \ifdim\@tempa\p@<\p@\scalebox{\@tempa}{\usebox\pandoc@box}%
  \else\usebox{\pandoc@box}%
  \fi%
}
% Set default figure placement to htbp
\def\fps@figure{htbp}
\makeatother



\ifLuaTeX
\usepackage[bidi=basic,shorthands=off]{babel}
\else
\usepackage[bidi=default,shorthands=off]{babel}
\fi
\ifLuaTeX
  \usepackage{selnolig} % disable illegal ligatures
\fi


\setlength{\emergencystretch}{3em} % prevent overfull lines

\providecommand{\tightlist}{%
  \setlength{\itemsep}{0pt}\setlength{\parskip}{0pt}}



 


\makeatletter
\@ifpackageloaded{caption}{}{\usepackage{caption}}
\AtBeginDocument{%
\ifdefined\contentsname
  \renewcommand*\contentsname{Table of contents}
\else
  \newcommand\contentsname{Table of contents}
\fi
\ifdefined\listfigurename
  \renewcommand*\listfigurename{List of Figures}
\else
  \newcommand\listfigurename{List of Figures}
\fi
\ifdefined\listtablename
  \renewcommand*\listtablename{List of Tables}
\else
  \newcommand\listtablename{List of Tables}
\fi
\ifdefined\figurename
  \renewcommand*\figurename{Figure}
\else
  \newcommand\figurename{Figure}
\fi
\ifdefined\tablename
  \renewcommand*\tablename{Table}
\else
  \newcommand\tablename{Table}
\fi
}
\@ifpackageloaded{float}{}{\usepackage{float}}
\floatstyle{ruled}
\@ifundefined{c@chapter}{\newfloat{codelisting}{h}{lop}}{\newfloat{codelisting}{h}{lop}[chapter]}
\floatname{codelisting}{Listing}
\newcommand*\listoflistings{\listof{codelisting}{List of Listings}}
\makeatother
\makeatletter
\makeatother
\makeatletter
\@ifpackageloaded{caption}{}{\usepackage{caption}}
\@ifpackageloaded{subcaption}{}{\usepackage{subcaption}}
\makeatother
\usepackage{bookmark}
\IfFileExists{xurl.sty}{\usepackage{xurl}}{} % add URL line breaks if available
\urlstyle{same}
% fallback for those not using the hyperref driver hyperxmp:
\makeatletter
\@ifundefined{xmpquote}{\newcommand{\xmpquote}[1]{#1}}{}
\makeatother
\hypersetup{
  pdftitle={Johnson \& Johnson Time Series Analysis},
  pdfauthor={Rubén A.},
  pdflang={en},
  colorlinks=true,
  linkcolor={blue},
  filecolor={Maroon},
  citecolor={Blue},
  urlcolor={Blue},
  pdfcreator={LaTeX via pandoc}}


\title{Johnson \& Johnson Time Series Analysis}
\usepackage{etoolbox}
\makeatletter
\providecommand{\subtitle}[1]{% add subtitle to \maketitle
  \apptocmd{\@title}{\par {\large #1 \par}}{}{}
}
\makeatother
\subtitle{A reproducible analysis with R and Quarto}
\author{Rubén A.}
\date{2026-10-16}
\begin{document}
\maketitle


\section{Introduction}\label{introduction}

The Johnson \& Johnson quarterly earnings per share series (1960--1980)
is a canonical example in the time series literature, appearing in
Shumway and Stoffer's \emph{Time Series Analysis and Its Applications}
as the first worked example of regression with seasonal dummies.

The series consists of 84 quarterly observations of earnings per share,
starting in 1960-Q1 and ending in 1980-Q4. It exhibits a strong upward
trend, clear quarterly seasonality, and cyclic structure at
business-cycle frequencies. These features make it an excellent case
study for comparing classical regression-based approaches with modern
time series models.

\subsection{Objective}\label{objective}

This report provides a reproducible analysis of the J\&J series using
three complementary approaches:

\begin{enumerate}
\def\labelenumi{\arabic{enumi}.}
\tightlist
\item
  \textbf{Regression baseline.} A log-linear model with a centered time
  trend and quarterly dummy variables, following the original exercise.
\item
  \textbf{Exploratory decomposition.} STL decomposition, seasonal plots,
  and quarterly boxplots to characterize the structure of the series.
\item
  \textbf{Alternative models and forecasting.} SARIMA and ETS models,
  validated via holdout and rolling cross-validation.
\end{enumerate}

The analysis is fully reproducible: all figures, tables, and numerical
results are generated from code in this repository.

\section{Data}\label{data}

The dataset used in this report is the \texttt{JohnsonJohnson} series
shipped with base R. It contains 84 quarterly observations, starting in
1960-Q1 and ending in 1980-Q4.

Earnings per share (EPS) is a standard financial metric computed by
dividing a company's net income by the number of outstanding shares. It
measures how much profit the company generated for each share of stock
during a given period. For Johnson \& Johnson, EPS is reported
quarterly, and this dataset tracks it from the first quarter of 1960
through the fourth quarter of 1980.

\begin{figure}

\centering{

\pandocbounded{\includegraphics[keepaspectratio]{report_files/figure-pdf/fig-original-1.pdf}}

}

\caption{\label{fig-original}Johnson \& Johnson quarterly earnings per
share, 1960--1980.}

\end{figure}%

As shown in Figure~\ref{fig-original}, the series displays two features
that shape the modeling strategy.

First, an upward trend: earnings per share roughly triple over the
twenty-one years.

Second, a clear seasonal oscillation: within each year, earnings rise
from the first to the third quarter and drop in the fourth. This pattern
repeats with remarkable consistency across the sample.

The variance also appears to grow with the level. Early years fluctuate
within a narrow band, while later years show much larger swings. This
pattern, called heteroscedasticity, motivates the log transformation
used throughout the analysis.

All subsequent work is performed on the log series, \(x_t = \log(y_t)\),
where \(y_t\) denotes the original earnings per share.

\section{Exploratory analysis}\label{exploratory-analysis}

Before fitting any model, we inspect the structure of the series. We use
the STL decomposition to separate trend and seasonality, quarterly
boxplots to check whether the seasonal pattern is stable across years,
and the sample ACF and PACF to guide model selection.

\subsection{STL decomposition}\label{stl-decomposition}

STL (Seasonal-Trend decomposition using LOESS) splits a series into
three additive components: a smooth trend, a repeating seasonal pattern,
and a remainder that captures everything else. It is a data-driven
alternative to fitting a regression first.

\begin{figure}

\centering{

\pandocbounded{\includegraphics[keepaspectratio]{output/stl.png}}

}

\caption{\label{fig-stl}STL decomposition of log earnings. Panels show
observed, trend, seasonal, and remainder components.}

\end{figure}%

Figure~\ref{fig-stl} shows four panels. The observed panel reproduces
the raw series, growing across the year and dropping in Q4, following
the same pattern seen in the boxplots. The trend panel is a smoothed
version of the series, with amplitude much smaller than the raw data.
The seasonal component is approximately periodic, oscillating between
about plus and minus one-tenth of a log unit. The remainder stays within
a small band around zero, roughly two-tenths of a unit on the log scale,
indicating that trend and seasonality capture most of the variation in
the series.

\subsection{Seasonal pattern}\label{seasonal-pattern}

A seasonal plot rearranges the series into a grid with one row per year
and one column per quarter. Each line traces the path of a single year
across the four quarters, and the color of the line encodes the year,
from grey (1960) to red (1980).

This layout makes the within-year pattern visible at a glance. If
earnings follow a stable seasonal shape, all lines rise and fall the
same way, and only their vertical position changes from year to year.
Deviations from the common shape stand out immediately.

\begin{figure}

\centering{

\pandocbounded{\includegraphics[keepaspectratio]{output/seasonal.png}}

}

\caption{\label{fig-seasonal}Seasonal plot of log earnings. Each line
represents one year; color encodes year from grey (1960) to red (1980).}

\end{figure}%

Figure~\ref{fig-seasonal} confirms a consistent within-year pattern:
earnings grow from Q1 to Q3 and decline in Q4. The color gradient shows
that later years (red) are systematically above earlier years (grey),
reflecting the upward trend. No year deviates visibly from this shape,
which justifies a fixed seasonal specification.

\subsection{Quarterly distribution}\label{quarterly-distribution}

A boxplot summarizes the distribution of a set of values in a compact
form. The central box spans the interquartile range, from the 25th to
the 75th percentile. The thick horizontal line inside the box marks the
median. The whiskers extend to the most extreme values within 1.5 times
the interquartile range, and points beyond them are outliers.

Here we group the 84 log earnings values by quarter. Each of the four
boxplots therefore summarizes 21 observations, one per year.

\begin{figure}

\centering{

\pandocbounded{\includegraphics[keepaspectratio]{output/quarter_boxplot.png}}

}

\caption{\label{fig-boxplot}Distribution of log earnings by quarter
across the 21 years.}

\end{figure}%

Figure~\ref{fig-boxplot} shows that the median increases from Q1 to Q3
and Q4 presents larger dispersion. Q3 exhibits the narrowest
interquartile range, indicating the least dispersion of log earnings in
that quarter. If we rearrange the quarters to place Q4 before Q1, a
pattern emerges that is easy to miss in the standard order: Q4 is
systematically above Q1, Q1 is the absolute minimum, and Q3 is the
absolute maximum. The within-year trajectory is therefore not monotonic
but cyclical, with a pronounced drop between Q3 and Q4 followed by a
further drop into Q1, and then a recovery through Q2 and Q3.

\subsection{ACF and PACF}\label{acf-and-pacf}

The sample autocorrelation function (ACF) measures the linear
correlation between the series and its own lagged values. The partial
autocorrelation function (PACF) does the same but removes the effect of
intermediate lags. Together they are the standard diagnostic for
choosing the orders of an ARIMA model.

\begin{figure}

\centering{

\pandocbounded{\includegraphics[keepaspectratio]{output/acf_pacf.png}}

}

\caption{\label{fig-acf}Time series (top), sample ACF (middle), and
sample PACF (bottom) of log earnings.}

\end{figure}%

Figure~\ref{fig-acf} shows that the ACF decays very slowly, which is the
classic signature of a non-stationary series. The PACF, by contrast, has
a single significant spike at lag 1. Beyond lag 1, it shows local minima
at lags 5, 9, and 13, all of which stay within the confidence band. The
isolated spike at lag 1 suggests a first-order autoregressive structure,
while the regular spacing of the local minima hints at a seasonal
component with a shift of one quarter. The slow decay of the ACF and the
shifted minima together motivate the seasonal differencing used in the
modeling section.

\section{Regression baseline}\label{regression-baseline}

The original exercise proposes a log-linear model with a centered time
trend and quarterly dummy variables:

\[
x_t = \beta \, t + \alpha_1 Q_1(t) + \alpha_2 Q_2(t) + \alpha_3 Q_3(t) + \alpha_4 Q_4(t) + \omega_t
\]

where \(x_t = \log(y_t)\), \(t\) is the time index centered at 1970,
\(Q_i(t) = 1\) if observation \(t\) belongs to quarter \(i\) and \(0\)
otherwise, and \(\omega_t\) is a Gaussian white noise sequence. The
model has no separate intercept because the four quarterly dummies span
the constant.

The parameters have a direct interpretation. The trend coefficient
\(\beta\) captures the average quarterly growth rate of log earnings
across the entire sample. Each \(\alpha_i\) represents the baseline
level of quarter \(i\) after removing the trend. The seasonal amplitude
is encoded in the differences \(\alpha_i - \alpha_j\) between quarters.

\subsection{Coefficient estimates}\label{coefficient-estimates}

The table below reports the estimated coefficients, their standard
errors, the test statistic, and the associated p-value for the
no-intercept specification. Each coefficient is tested against zero: a
small p-value indicates that the coefficient is statistically different
from zero.

{

\begin{longtable}[]{@{}lrrrr@{}}

\caption{\label{tbl-coef}Coefficient estimates for the no-intercept
regression.}

\tabularnewline

\toprule\noalign{}
term & estimate & std.error & statistic & p.value \\
\midrule\noalign{}
\endhead
\bottomrule\noalign{}
\endlastfoot
trend & 0.1672 & 0.0023 & 73.9990 & 0 \\
quarterQ1 & 1.0528 & 0.0274 & 38.4804 & 0 \\
quarterQ2 & 1.0809 & 0.0274 & 39.4999 & 0 \\
quarterQ3 & 1.1510 & 0.0274 & 42.0350 & 0 \\
quarterQ4 & 0.8823 & 0.0274 & 32.1858 & 0 \\

\end{longtable}

}

Table~\ref{tbl-coef} shows the estimated coefficients. The trend
estimate \(\hat{\beta}\) is positive and highly significant, confirming
the upward movement visible in the series. All four quarterly
coefficients are also significant at conventional levels. They differ
from each other, and their spread reflects the amplitude of the seasonal
pattern identified in the exploratory analysis.

\subsection{Observed vs fitted}\label{observed-vs-fitted}

Overlaying fitted values on the observed series is the quickest way to
judge how well a model tracks the data. Two failure modes are common: a
model that captures the trend but misses the seasonality, or one that
captures the seasonality but misses long-run movements.

\begin{figure}

\centering{

\pandocbounded{\includegraphics[keepaspectratio]{output/observed_vs_fitted.png}}

}

\caption{\label{fig-observed}Observed log earnings (black) and fitted
values from the no-intercept regression (red).}

\end{figure}%

Figure~\ref{fig-observed} shows that the regression captures the overall
upward trend and the quarterly pattern, but the fitted values form a
smooth quasi-periodic curve that misses the year-to-year variation in
the observed series. This is expected: a regression with five parameters
cannot capture business-cycle dynamics. The residuals, the difference
between the two curves, are examined next.

\subsection{Residual diagnostics}\label{residual-diagnostics}

Residuals are the differences between observed and fitted values. If the
model captures all systematic structure, the residuals should behave
like white noise: no trend, no seasonal pattern, no autocorrelation.
Non-random residuals signal that something is missing from the
specification.

\begin{figure}

\centering{

\pandocbounded{\includegraphics[keepaspectratio]{output/residuals.png}}

}

\caption{\label{fig-residuals}Residuals from the no-intercept
regression.}

\end{figure}%

The residuals in Figure~\ref{fig-residuals} are not white noise. They
display a low-frequency oscillation with a clear structure: a sustained
negative phase through the 1960s, a sharp positive peak centered around
1972, and a return to negative values after 1976. The peak is robust
across moving-average windows of different lengths, which confirms it is
a genuine feature of the residuals rather than an artifact of smoothing.
The oscillation is not strictly periodic and its duration is not fixed,
which is characteristic of economic cycles rather than of a mechanical
seasonal pattern. The regression specification, with only five
parameters and a rigid trend, cannot represent this structure,
motivating the use of autoregressive models such as SARIMA in the next
section.

\section{Stationarity}\label{stationarity}

Most time series models, including SARIMA and ETS, assume the series is
stationary after appropriate transformations. A series is weakly
stationary if its mean, variance, and autocovariance structure do not
change over time. We test this assumption with two complementary
stationarity tests, and we also apply the Ljung-Box test to check
whether the series exhibits serial dependence.

\subsection{ADF and KPSS}\label{adf-and-kpss}

The Augmented Dickey-Fuller (ADF) test has the null hypothesis that the
series contains a unit root, meaning it is non-stationary. The
Kwiatkowski-Phillips-Schmidt-Shin (KPSS) test has the opposite null:
that the series is stationary. Using both tests together reduces the
risk of a misleading conclusion.

{

\begin{longtable}[]{@{}
  >{\raggedright\arraybackslash}p{(\linewidth - 10\tabcolsep) * \real{0.2073}}
  >{\raggedright\arraybackslash}p{(\linewidth - 10\tabcolsep) * \real{0.0610}}
  >{\raggedleft\arraybackslash}p{(\linewidth - 10\tabcolsep) * \real{0.1220}}
  >{\raggedleft\arraybackslash}p{(\linewidth - 10\tabcolsep) * \real{0.0976}}
  >{\raggedright\arraybackslash}p{(\linewidth - 10\tabcolsep) * \real{0.3293}}
  >{\raggedright\arraybackslash}p{(\linewidth - 10\tabcolsep) * \real{0.1829}}@{}}

\caption{\label{tbl-stationarity}Stationarity tests on the log series
and its first difference.}

\tabularnewline

\toprule\noalign{}
\begin{minipage}[b]{\linewidth}\raggedright
series
\end{minipage} & \begin{minipage}[b]{\linewidth}\raggedright
test
\end{minipage} & \begin{minipage}[b]{\linewidth}\raggedleft
statistic
\end{minipage} & \begin{minipage}[b]{\linewidth}\raggedleft
p\_value
\end{minipage} & \begin{minipage}[b]{\linewidth}\raggedright
null\_hypothesis
\end{minipage} & \begin{minipage}[b]{\linewidth}\raggedright
conclusion
\end{minipage} \\
\midrule\noalign{}
\endhead
\bottomrule\noalign{}
\endlastfoot
log series & ADF & -1.1543 & 0.9087 & Unit root (non-stationary) &
Non-stationary \\
log series & KPSS & 2.1867 & 0.0100 & Level-stationary &
Non-stationary \\
first difference & ADF & -4.5649 & 0.0100 & Unit root (non-stationary) &
Stationary \\
first difference & KPSS & 0.1720 & 0.1000 & Level-stationary &
Stationary \\

\end{longtable}

}

Table~\ref{tbl-stationarity} shows that both tests agree on the log
series: ADF fails to reject the unit-root null, and KPSS rejects
stationarity. On the first difference, both tests agree in the opposite
direction. The conclusion is that the log series is non-stationary and a
single regular difference is enough to make it stationary.

\subsection{Ljung-Box test}\label{ljung-box-test}

The Ljung-Box test checks whether a series exhibits significant
autocorrelation up to a given lag. Applied to the log series, it
confirms the visual impression from the ACF.

{

\begin{longtable}[]{@{}rrr@{}}

\caption{\label{tbl-ljungbox}Ljung-Box test on the log series, lag =
10.}

\tabularnewline

\toprule\noalign{}
statistic & df & p\_value \\
\midrule\noalign{}
\endhead
\bottomrule\noalign{}
\endlastfoot
613.0387 & 10 & 0 \\

\end{longtable}

}

The Ljung-Box test checks whether a series shows autocorrelation up to a
given lag. Its null hypothesis is that the values are independent over
time. The test statistic of 613 with 10 degrees of freedom produces a
p-value below any conventional threshold. The null is rejected
decisively: the log series has strong serial dependence, and its values
are far from independent. This provides one more reason to move beyond a
regression with independent errors toward a model that explicitly
captures the autocorrelation, such as SARIMA.

\section{Alternative models}\label{alternative-models}

The regression baseline captures trend and seasonality but leaves a
cyclic residual structure unexplained. We now fit two families of models
specifically designed to handle autocorrelation: SARIMA and ETS.

\subsection{Model comparison}\label{model-comparison}

How should we choose between SARIMA and ETS? Both model the same series,
but they represent different assumptions about the underlying structure.
Information criteria provide a principled comparison. AIC and BIC trade
off goodness of fit against model complexity: lower values indicate a
better balance. They are not directly comparable across different
datasets, but they are meaningful when comparing models fitted to the
same data.

{

\begin{longtable}[]{@{}lrrr@{}}

\caption{\label{tbl-model-comparison}Information criteria for the SARIMA
and ETS fits on the full log series.}

\tabularnewline

\toprule\noalign{}
model & log\_lik & aic & bic \\
\midrule\noalign{}
\endhead
\bottomrule\noalign{}
\endlastfoot
SARIMA & 82.47 & -154.95 & -143.04 \\
ETS & 19.28 & -20.56 & 1.32 \\

\end{longtable}

}

Table~\ref{tbl-model-comparison} compares the two models using
log-likelihood, AIC, and BIC. The SARIMA model has an AIC of -154.95,
about 134.4 points lower than ETS. The BIC difference follows the same
direction and is even larger in magnitude. All three criteria therefore
point to the SARIMA model as the better trade-off between fit and
complexity.

\subsection{SARIMA specification}\label{sarima-specification}

SARIMA extends ARIMA to handle seasonal structure. Its full
specification is \(\text{ARIMA}(p, d, q)(P, D, Q)[s]\), where:

\begin{itemize}
\tightlist
\item
  \(p, d, q\) are the regular autoregressive, differencing, and
  moving-average orders.
\item
  \(P, D, Q\) are the seasonal counterparts.
\item
  \(s\) is the seasonal period, equal to 4 for quarterly data.
\end{itemize}

We fit the model automatically with \texttt{forecast::auto.arima}, which
searches over a grid of orders and selects the one minimizing AICc.

The search selects a SARIMA(2,0,0)(1,1,0){[}4{]} with drift. The two
regular autoregressive terms capture short-run dependence between
consecutive quarters. The seasonal term at lag 4 captures the annual
cycle. The drift term reflects the average quarterly growth of log
earnings. Because the seasonal differencing order is \(D = 1\), the
model effectively works with year-over-year changes, which removes the
seasonal non-stationarity detected earlier.

\subsection{ETS specification}\label{ets-specification}

ETS stands for \emph{Error, Trend, Seasonal}. Each component can take
different forms: error can be additive (A) or multiplicative (M); trend
can be none (N), additive (A), or damped (Ad); seasonality can be none
(N), additive (A), or multiplicative (M). The combination defines a
family of exponential smoothing models.

We fit the model automatically with \texttt{forecast::ets}, which
evaluates all feasible combinations and selects the one minimizing AICc.

The selected model is ETS(A,A,A), also known as Holt-Winters additive.
It decomposes the series into level, trend, and seasonality, each
updated exponentially with weights \(\alpha\), \(\beta\), and
\(\gamma\). The estimated seasonal smoothing parameter is 0.54,
indicating that the seasonal component adapts substantially each year.

\section{Forecast validation}\label{forecast-validation}

Comparing models on the same data used for fitting can be misleading. We
evaluate out-of-sample performance with two complementary strategies: a
single holdout split and rolling cross-validation.

\subsection{Holdout validation}\label{holdout-validation}

The series is split into a training set covering 1960 to 1975 (64
observations) and a test set covering 1976 to 1980 (20 observations).
Both models are refit on the training set and evaluated on the test set.

Three metrics summarize the forecast errors:

\begin{itemize}
\tightlist
\item
  \textbf{RMSE} (Root Mean Squared Error) is the square root of the
  average squared error. It penalizes large errors more heavily than
  small ones and is on the same scale as the log earnings.
\item
  \textbf{MAE} (Mean Absolute Error) is the average absolute error. It
  is less sensitive to outliers than RMSE.
\item
  \textbf{MAPE} (Mean Absolute Percentage Error) expresses the error as
  a percentage of the observed value. It is useful for comparing across
  scales, but it is unstable when the true value is close to zero.
\end{itemize}

{

\begin{longtable}[]{@{}lrrr@{}}

\caption{\label{tbl-holdout}Holdout accuracy for SARIMA and ETS on the
1976--1980 test set.}

\tabularnewline

\toprule\noalign{}
model & RMSE & MAE & MAPE \\
\midrule\noalign{}
\endhead
\bottomrule\noalign{}
\endlastfoot
SARIMA & 0.0833 & 0.0712 & 2.9650 \\
ETS & 0.2181 & 0.1815 & 7.1837 \\

\end{longtable}

}

Table~\ref{tbl-holdout} shows that SARIMA achieves a test RMSE of
0.0833, a MAE of 0.0712, and a MAPE of 2.96 percent. ETS, in contrast,
achieves an RMSE of 0.2181, a MAE of 0.1815, and a MAPE of 7.18 percent.
The difference is substantial: SARIMA is roughly 2.6 times more accurate
on the test set by RMSE.

\subsection{Rolling cross-validation}\label{rolling-cross-validation}

Rolling cross-validation fits the model at each time point and predicts
one step ahead, then averages the squared errors across all folds. This
uses every observation as a potential test point.

{

\begin{longtable}[]{@{}lr@{}}

\caption{\label{tbl-tscv}One-step-ahead rolling cross-validation RMSE.}

\tabularnewline

\toprule\noalign{}
model & rmse \\
\midrule\noalign{}
\endhead
\bottomrule\noalign{}
\endlastfoot
SARIMA & 0.1300 \\
ETS & 0.1389 \\

\end{longtable}

}

Table~\ref{tbl-tscv} shows the same ranking: SARIMA has a lower
cross-validated RMSE than ETS. Both metrics agree, giving confidence in
the conclusion.

\section{Discussion}\label{discussion}

Three findings stand out from the analysis.

First, the log transformation and the choice of a single difference are
well justified by the data. The raw series shows variance growing with
the level, and the log scale stabilizes it. The STL decomposition then
reveals that trend and seasonality account for most of the variation,
leaving a remainder with amplitude of roughly ±0.2 on the log scale. The
ADF and KPSS tests on the differenced series agree that a single regular
difference suffices to achieve stationarity. This provides a clear
justification for the \texttt{d\ =\ 1} setting that the modeling section
assumes.

Second, the regression baseline is useful as an interpretable model but
insufficient as a complete description of the series. All five
coefficients are significant, and the fitted values track the trend and
the quarterly pattern accurately. However, the residuals are not white
noise: they display a low-frequency oscillation with irregular phases,
including a sustained negative period through the 1960s, a sharp peak
around 1972, and a return to negative values after 1976. The
five-parameter specification cannot represent this structure. This is a
genuine limitation of the model class, not a numerical artifact.

Third, the SARIMA model outperforms ETS on every metric considered. The
automatic search selects a SARIMA(2,0,0)(1,1,0){[}4{]} with drift, which
combines regular autoregressive terms, a seasonal autoregressive term at
lag 4, and seasonal differencing. The ETS model selects the additive
Holt-Winters specification, ETS(A,A,A). SARIMA achieves an AIC of about
134 points lower than ETS, and a test-set RMSE that is roughly 2.6 times
smaller. The rolling cross-validation confirms the same ranking. The
most likely explanation is that the regular AR terms in SARIMA capture
the short-run autocorrelation that the additive ETS specification cannot
represent, while the seasonal AR term captures the annual cycle.

Two caveats apply to the third finding. First, the seasonal differencing
order (\texttt{D\ =\ 1}) chosen by \texttt{auto.arima} differs from the
regular differencing order (\texttt{d\ =\ 1}) suggested by the ADF and
KPSS tests. Both approaches remove the non-stationarity, but they do so
through different mechanisms. The AICc criterion preferred the
seasonal-differencing solution for this series. Second, the holdout RMSE
and the tsCV RMSE measure different things: the holdout evaluates a
single model trained on 64 observations and asked to forecast 20 steps
ahead, while the tsCV fits a new model at every step and predicts one
step ahead. Both metrics agree on the ranking of the two models, which
strengthens the conclusion.

\section{Limitations}\label{limitations}

The analysis has several limitations worth noting.

The sample is small: 84 quarterly observations spanning 1960 to 1980.
With only 21 years of data, the effective sample size for estimating
seasonal and business-cycle structure is limited. Forecasts beyond 1980
should not be extrapolated without additional data, since the later
behavior of the series may reflect economic conditions absent from this
sample.

The comparison between SARIMA and ETS relies on automatic model
selection. \texttt{auto.arima} and \texttt{ets} search over candidate
specifications by AICc, and the selected models are the best within
their respective families, but not necessarily the best possible models
for this series. A manual search guided by the ACF and PACF could yield
a different SARIMA specification. Similarly, the ranking of the two
families could change under a different loss function.

The SARIMA and ETS models were compared on a single train/test split
(1960--1975 for training, 1976--1980 for testing), and this split was
also supplemented by rolling cross-validation. The single split places
the test period in the last five years of the sample, when the series
appears to flatten slightly. A different split, or a larger number of
folds in the cross-validation, could change the magnitude of the
reported differences, although the direction of the ranking was
consistent across both evaluations.

The seasonal plot shows no year that deviates strongly from the overall
shape, but this does not rule out structural breaks at a finer scale. A
formal breakpoint analysis, such as the Chow test or the Bai--Perron
procedure, was not performed. If the series contains a structural break,
the models presented here would fit neither regime well.

Finally, all figures and tables in this report are generated from a
single pipeline. The pipeline is reproducible, but it assumes that the
input series is \texttt{datasets::JohnsonJohnson} unchanged. Modifying
the input or the transformation would require rerunning the analysis and
updating the interpretation accordingly.

\section{Reproducibility}\label{reproducibility}

All results in this report are generated from code in this repository.
The analysis is organized into two stages.

The first stage runs the full pipeline: data loading, exploratory
decomposition, regression fitting, stationarity tests, alternative
models, and forecast validation. It writes all figures and tables to
\texttt{output/} and saves the key R objects to
\texttt{output/results.rds}.

Rscript analysis.R On systems where Rscript is not on the PATH, the same
command can be run from within an R session opened at the project root:
source(``analysis.R'') The second stage renders this report from the
precomputed results. It does not refit any model; it reads
output/results.rds and produces the final document. quarto render
report.qmd To regenerate the PDF version with the LaTeX source preserved
for external editing: quarto render report.qmd --to pdf The LaTeX source
is written to report.tex alongside the PDF.

Package versions are pinned via renv. To restore the exact environment
used to produce this report: renv::restore() The R session used to
compile the report is recorded below. R version: R version 4.4.1
(2024-06-14 ucrt)

Key package versions:

\begin{verbatim}
Registered S3 method overwritten by 'quantmod':
  method            from
  as.zoo.data.frame zoo 
\end{verbatim}

forecast : 8.23.0 tseries : 0.10.58 ggplot2 : 3.5.1 broom : 1.0.7 knitr
: 1.48 dplyr : 1.1.4 tibble : 3.2.1 rlang : 1.1.4 quarto : 1.5.1




\end{document}
